% DRAFT 2026-10-28 -- replacement for Thesis_Optimization_Background.tex,
% \subsection{Architecture E: Segmented Multi-Law Guidance} (lines 589-636 at Overleaf 9bbba1d).
% Not applied. Figures: dev-notes/figures/fig_segmented_waypoints.py (new),
% fig_segmented_schedule.py (redrawn), fig_thesis_arc_structures.py (row E).

\subsection{Architecture E: Segmented Multi-Law Guidance}
\label{ssec:segmented}

Architectures A to D fly one guidance law from second-stage ignition to cut-off. Architecture E
divides the ascent among several laws, so that each part of it can be flown by the law suited to
it, a passive law through the dense atmosphere and closed-loop guidance above it
(Section~\ref{sec:guidance_classification}). It also gives the burn before the coast a target of
its own, where architecture B aims it at the final orbit.

\paragraph{Schedule.}
A schedule is an ordered list of $n$ laws with activation altitudes $0 = h_1 < h_2 < \dots < h_n$.
Law $k$ steers from the instant the vehicle climbs through $h_k$ until it climbs through
$h_{k+1}$: the first takes over at the pitch-over kick and the last flies to insertion. A hand-off
changes the steering objective, not the vehicle, so it is independent of staging and may fall in
the first stage. The segment laws are the gravity turn, the Apollo law, vector PEG, the two tangent
laws, and the indirect law, which replays the reference's optimal $\alpha$ at the current altitude
instead of running live. Any of them may occupy any position; Table~\ref{tab:segment_roles} lists
what each aims at and the difficulty it meets in each position.

\begin{table}[htbp]
  \centering
  \caption{The segment laws: what each aims at, and the difficulty it meets as the first,
    an intermediate or the final law.}
  \label{tab:segment_roles}
  \begin{tabular}{p{0.13\linewidth} p{0.15\linewidth} p{0.19\linewidth} p{0.19\linewidth} p{0.22\linewidth}}
    \toprule
    \textbf{Law} & \textbf{Aims at} & \textbf{First} & \textbf{Intermediate} & \textbf{Final} \\
    \midrule
    Gravity turn & nothing, $\alpha = 0$ & none & cannot reach its waypoint; the next law inherits the miss & insertion set by the kick and burn timing alone \\
    Apollo & full state at a set time & steers at high dynamic pressure, where lift is neglected & none & burn lengths must be set externally \\
    Vector PEG & full state, own $t_{go}$ & steers at high dynamic pressure & a coefficient freeze below about \qty{10}{\s} destabilizes its corrector & a burn shorter than the freeze flies one frozen solution \\
    Tangent laws & flight-path angle only & steers at high dynamic pressure & altitude and speed at the waypoint uncontrolled & radius and speed set by burn timing alone \\
    Indirect & stored optimal $\alpha$ & replays the reference's gravity turn & open loop: departures persist & open loop: insertion errors uncorrected \\
    \bottomrule
  \end{tabular}
\end{table}

\paragraph{Targets and timing.}
Law $k < n$ aims at the waypoint $(h_{k+1}, v_{k+1}, \gamma_{k+1})$, the state of the indirect
reference of Section~\ref{ssec:indirect_pmp_impl} where it climbs through $h_{k+1}$, and the final
law at the target orbit (Figure~\ref{fig:segmented_waypoints}). The reference's optimal control
acts in the second stage only, so below second-stage ignition its waypoints are states of a
gravity turn. The time-to-go of Eq.~\eqref{eq:tgo_rocket} collapses at the stage boundary, so a
law that needs one counts down instead to the time the reference takes to the waypoint,
\begin{equation}
  t_{go} = t_{\mathrm{deadline}} - t, \qquad
  t_{\mathrm{deadline}} = t_0 + t_{\mathrm{ref}}(h_{k+1}) - t_{\mathrm{ref}}\bigl(h(t_0)\bigr),
  \label{eq:segment_deadline}
\end{equation}
where $t_0$ is the instant law $k$ takes over or a burn restarts, and $t_{\mathrm{ref}}(h)$ is the
reference's time at altitude $h$ on its ascent.

\begin{figure}[htbp]
  \centering
  \includegraphics[width=\linewidth]{Figures/segmented_waypoints.png}
  \caption{Targets of a three-law schedule. Each law aims at the reference's state at the next
    hand-off altitude, and the last at the target orbit.}
  \label{fig:segmented_waypoints}
\end{figure}

\paragraph{Ending the second-stage burns.}
The second stage is flown thrust--coast--thrust, as in architecture B, and the hand-offs fall in it
wherever their altitudes place them. Two rules end its burns (Table~\ref{tab:segmented_forms}). In
the swarm-timed form the swarm sets both burn lengths and the final law aims at the orbit in both,
so the design vector is that of Eq.~\eqref{eq:pso_coast_vector} extended by the hand-off altitudes,
\begin{equation}
  \mathbf{x} = \bigl[\,\Delta t_c,\; \Delta t_r\%,\; \text{coast-start}\%,\; \gamma_p,\;
                     f_2, \dots, f_n\,\bigr].
  \label{eq:segmented_vector_swarm}
\end{equation}
In the law-terminated form the final law ends each burn on its own time-to-go. It aims the first
burn at the \emph{coast-start waypoint}, the reference's state where its first second-stage burn
ends, coasts for $\Delta t_c$, and aims the second at the orbit. The burn lengths become outputs,
\begin{equation}
  \mathbf{x} = \bigl[\,\Delta t_c,\; \gamma_p,\; f_2, \dots, f_n\,\bigr].
  \label{eq:segmented_vector_lt}
\end{equation}
This form needs a final law with a cut-off rule of its own, which among the implemented laws only
vector PEG has, and every hand-off must precede the coast-start waypoint. It is the form flown in
Chapter~\ref{chapter:results_and_discussion} (Figure~\ref{fig:segmented_schedule}).

\begin{table}[htbp]
  \centering
  \caption{The two forms of architecture E.}
  \label{tab:segmented_forms}
  \begin{tabular}{p{0.24\linewidth} p{0.33\linewidth} p{0.33\linewidth}}
    \toprule
     & \textbf{Swarm-timed} & \textbf{Law-terminated} \\
    \midrule
    Burns ended by & the swarm & the final law's time-to-go \\
    First burn aims at & the target orbit & the coast-start waypoint \\
    Final law & any & one with its own cut-off rule \\
    Hand-offs & up to the target orbit & below the coast-start waypoint \\
    Design vector & Eq.~\eqref{eq:segmented_vector_swarm} & Eq.~\eqref{eq:segmented_vector_lt} \\
    \bottomrule
  \end{tabular}
\end{table}

\begin{figure}[htbp]
  \centering
  \includegraphics[width=\linewidth]{Figures/segmented_schedule.png}
  \caption{A law-terminated schedule: a gravity turn hands over to vector PEG, which ends both
    burns on its own time-to-go. The swarm chooses the kick and the coast.}
  \label{fig:segmented_schedule}
\end{figure}

\paragraph{Hand-off altitudes.}
The hand-off altitudes are fixed or searched with the other variables. Searched, they enter
$\mathbf{x}$ as fractions $f_k \in [0, 1]$ mapped cumulatively onto
$[h_{\mathrm{lb}}, h_{\mathrm{ub}}]$,
\begin{equation}
  h_2 = h_{\mathrm{lb}} + f_2\,(h_{\mathrm{ub}} - h_{\mathrm{lb}}), \qquad
  h_k = h_{k-1} + f_k\,(h_{\mathrm{ub}} - h_{k-1}), \quad k = 3, \dots, n,
  \label{eq:handoff_mapping}
\end{equation}
so that every particle yields a strictly increasing schedule and none is rejected for ordering the
laws out of sequence. Both forms minimize the penalized burn-time sum $J$ of
Section~\ref{ssec:pso}, as architecture B does.

% ---------------------------------------------------------------------------
% Knock-on edit, same file, Table tab:architectures, row E:
%   OLD: E\; Segmented multi-law & several laws, coast, thrust & those of B and the law activation altitudes & particle swarm \\
%   NEW: E\; Segmented multi-law & thrust, coast, thrust, laws handed over by altitude & pitch-over angle, coast and hand-off altitudes; B's burn timing when swarm-timed & particle swarm \\
